\section{引言}
遮挡交互操作要求机器人根据动态可见性、对象间关系及任务进展，
持续决定合适的运动。模仿学习可以从示范中学习这些决策，但其常见优化目标是
预测示范指令，而评测目标是整个任务是否完成。少量关键时刻的错误足以造成拨空、
推离目标或未完成复原；大量普通运动帧上的准确预测可能掩盖这些错误。

目前实机已经观察到从基础ACT到语义表示和历史条件纠偏的逐步收益，
详见第\ref{sec:evidence}节。研究不是仅凭直觉断言模块有效，
而是从这些证据提出机制假设，再用理论明确条件并设计可反驳的对照。
因此，本文的问题不是如何预先规定快慢网络，而是：
\emph{在有限示范与受限在线计算下，如何使学到的策略更可靠地依赖任务相关状态，
并在执行过程中利用新证据修正偏差？}
我们遵循三个层次回答。第一，保留RGB可见细节，同时用语义分割提供对象和空间范围的
显式表达，降低单纯从动作标签中发现相关视觉结构的负担。
第二，用任务几何监督查询token，要求策略内部表示保留可见程度、目标位置及工具关系，
而不只满足全局动作拟合。第三，在基础动作块执行期间，用最新观测、
当前计划和短期历史学习目标时间对齐的局部残差。

这条路径不依赖以下错误前提：更多输入一定有利，较低训练误差一定提高成功率，
或者高频重规划本身一定优于持续执行。
行为克隆的序列分布偏移已有明确分析\cite{ross2011}；
因果混淆研究也表明，额外观测及历史可能强化错误相关性\cite{dehaan2019}。
ACT的时间集成本身包含频繁查询\cite{zhao2023act}，不能视为完全没有反馈的基线。
我们的分析需要解释的是表示、历史和反馈的\emph{使用方式}，不是简单比较模块数量。

\section{相关理论与分析定位}
性能差分恒等式把策略性能变化写为新策略访问分布下的优势累积，
例如文献\cite{schulman2015}的式(1)；本文给出终端失败事件下的有限时域版本，
不继承其优化算法的单调改进保证。
DAgger研究训练分布与策略访问分布之间的不一致\cite{ross2011}。
其中经典最坏情况误差累积结论采用相应分类损失与代价假设，
不能直接把ACT的L1数值代入其公式。
稳定性约束模仿学习进一步说明动力学的增量稳定性质影响样本复杂度和轨迹偏差
\cite{tu2022}；本文没有学习或验证这样的全局稳定性证书。

本文后续的平方风险分解属于条件期望投影分析。
它解释理想信息收益、表示损失及有限训练误差如何竞争，但本身不是样本复杂度改进定理。
残差动作块修正已有相关研究\cite{a2c2}，故本文不将一般实时残差形式声明为首创；
需要验证的是低人工成本语义监督、几何查询表示与历史条件纠偏在目标任务上的组合价值。
